% ===== PRÉAMBULE CANONIQUE — à coller VERBATIM en tête de CHAQUE .tex =====
\documentclass[11pt,a4paper]{article}
\usepackage[utf8]{inputenc}\usepackage[T1]{fontenc}\usepackage{lmodern}
\usepackage[french]{babel}
\usepackage{amsmath,amssymb,mathtools}\usepackage{icomma}
\usepackage[version=4]{mhchem}          % \ce{...} : équations & formules
\usepackage{siunitx}\sisetup{locale=FR} % \qty{}{}, \si{}, \ang{}
\usepackage{chemfig}                    % \chemfig{...} : structures développées
\usepackage{xcolor}\usepackage{colortbl}
\usepackage{tikz}\usepackage{pgfplots}\pgfplotsset{compat=1.18}
\usepackage[most]{tcolorbox}\usepackage{enumitem}\usepackage{array,booktabs}
\usepackage[a4paper,margin=2.2cm]{geometry}
\usetikzlibrary{arrows.meta,calc,angles,quotes,positioning,patterns,backgrounds,shapes.geometric}
\definecolor{primary}{RGB}{222,49,99}\definecolor{primarydark}{RGB}{150,22,55}
\definecolor{defcol}{RGB}{0,114,178}\definecolor{methcol}{RGB}{52,92,148}
\definecolor{excol}{RGB}{0,135,90}\definecolor{attncol}{RGB}{200,40,55}
\tcbset{boxbase/.style={enhanced,breakable,boxrule=0.9pt,arc=2.5pt,
  left=3mm,right=3mm,top=2.3mm,bottom=2.3mm,fonttitle=\bfseries,coltitle=white}}
% --- boîtes TUTORIEL (noms EXACTS, ne pas renommer) ---
\newtcolorbox{cle}[1][]{boxbase,colback=defcol!6,colframe=defcol,
  title={\textsc{À retenir}\ifx\relax#1\relax\relax\else\textmd{ : \itshape #1}\fi}}
\newtcolorbox{methode}[1][]{boxbase,colback=methcol!7,colframe=methcol,sharp corners,
  title={\textsc{Méthode}\ifx\relax#1\relax\relax\else\textmd{ --- \itshape #1}\fi}}
\newtcolorbox{exemplebox}[1][]{boxbase,colback=excol!5,colframe=excol,
  title={\textsc{Exemple}\ifx\relax#1\relax\relax\else\textmd{ : \itshape #1}\fi}}
\newtcolorbox{piege}[1][]{boxbase,colback=attncol!6,colframe=attncol,
  title={\textsc{Piège}\ifx\relax#1\relax\relax\else\textmd{ : \itshape #1}\fi}}
\newtcolorbox{remarque}[1][]{boxbase,colback=black!4,colframe=black!45,
  title={\textsc{Remarque}\ifx\relax#1\relax\relax\else\textmd{ : \itshape #1}\fi}}
% --- boîtes EXERCICE / CORRIGÉ (noms EXACTS, auto-comptées) ---
\newtcolorbox[auto counter]{exo}[1][]{boxbase,colback=primary!4,colframe=primary,
  title={\textsc{Exercice}~\thetcbcounter\ifx\relax#1\relax\relax\else\textmd{ --- \itshape #1}\fi}}
\newtcolorbox[auto counter]{corrige}[1][]{boxbase,colback=excol!4,colframe=excol,
  title={\textsc{Corrigé}~\thetcbcounter\ifx\relax#1\relax\relax\else\textmd{ --- \itshape #1}\fi}}
\newcommand{\R}{\mathbb{R}}\newcommand{\N}{\mathbb{N}}\newcommand{\Z}{\mathbb{Z}}\newcommand{\ds}{\displaystyle}
% (un \usepackage{fancyhdr}+en-tête est possible ; il est ignoré côté web)
% =========================================================================
\begin{document}
\begin{center}\color{primarydark}\large\bfseries Thème 4 — Corrigé \quad$\bullet$\quad Niveau Moyen\end{center}

\begin{corrige}[l'effet de la dilution]
$E=0{,}34+\dfrac{0{,}06}{2}\log[\ce{Cu^{2+}}]=0{,}34+0{,}03\log[\ce{Cu^{2+}}]$.
\begin{enumerate}[label=\alph*)]
  \item $E=0{,}34+0{,}03\times 0=+0{,}34$~V.
  \item $E=0{,}34+0{,}03\times(-1)=+0{,}31$~V.
  \item $E=0{,}34+0{,}03\times(-2)=+0{,}28$~V.
\end{enumerate}
Plus on \textbf{dilue} l'oxydant \ce{Cu^{2+}}, plus le potentiel \textbf{diminue} : la demi-pile
devient moins oxydante.
\end{corrige}

\begin{corrige}[une demi-pile d'argent]
$n=1$, $[\mathrm{Red}]=1$ (métal) :
\[
E=0{,}80+0{,}06\log(0{,}10)=0{,}80+0{,}06\times(-1)=+0{,}74~\text{V}.
\]
\end{corrige}

\begin{corrige}[f.é.m.\ hors des conditions standard]
\begin{enumerate}[label=\alph*)]
  \item Cathode (argent, $n=1$) : $E_{\text{cath}}=0{,}80+0{,}06\log(0{,}10)=+0{,}74$~V.
  \item Anode (cuivre, $n=2$) : $E_{\text{anode}}=0{,}34+0{,}03\log(1)=+0{,}34$~V.
  \item $\Delta E=E_{\text{cath}}-E_{\text{anode}}=0{,}74-0{,}34=+0{,}40$~V.
\end{enumerate}
\end{corrige}

\begin{corrige}[le potentiel du permanganate]
Le rapport vaut $\dfrac{0{,}01\times 1^{8}}{0{,}01}=1$, donc $\log(1)=0$ :
\[
E=1{,}51+\frac{0{,}06}{5}\times 0=+1{,}51~\text{V}.
\]
À $\mathrm{pH}=0$ et concentrations égales, on retrouve exactement $E_{\mathrm{r}}^{\circ}$.
\end{corrige}

\begin{corrige}[retrouver un rapport de concentrations]
$0{,}83=0{,}77+0{,}06\log\dfrac{[\ce{Fe^{3+}}]}{[\ce{Fe^{2+}}]}$, donc
$0{,}06\log(r)=0{,}06\Rightarrow\log(r)=1\Rightarrow r=10$. Le rapport
$\dfrac{[\ce{Fe^{3+}}]}{[\ce{Fe^{2+}}]}=10$.
\end{corrige}

\end{document}
